\subsubsection*{Question 1}
On écrit $5=\dfrac{10}{2}$. Ainsi :
\[5-\dfrac{3}{2}=\dfrac{10}{2}-\dfrac{3}{2}=\dfrac{7}{2}.\]

La bonne réponse est \textbf{b.}

\subsubsection*{Question 2}
On utilise la propriété $(a^m)^n=a^{m \times n}$ :
\[\left(2^3\right)^{-5}=2^{3\times(-5)}=2^{-15}.\]

La bonne réponse est \textbf{b.}

\subsubsection*{Question 3}
Une baisse de $20\,\%$ correspond à un coefficient multiplicateur égal à : $1-\dfrac{20}{100}=0{,}8$.

Le nouveau prix est donc $990\times 0{,}8$.

La bonne réponse est \textbf{d.}

\subsubsection*{Question 4}
Obtenir un nombre supérieur ou égal à 5 signifie obtenir 5 ou 6. Donc, dans le cas présent :
\[P(5 \cup 6)=P(5)+P(6)=0{,}3+0{,}3=0{,}6.\]

La bonne réponse est \textbf{c.}

\subsubsection*{Question 5}
Un produit est nul si, et seulement si, l'un de ses facteurs est nul :
\[x-2=0 \quad \text{ou} \quad 2x+1=0 \Rightarrow 2x=-1 \Rightarrow x=-\dfrac{1}{2}.\]
On obtient donc $x=2$ ou $x=-\dfrac{1}{2}$.

La bonne réponse est \textbf{a.}

\subsubsection*{Question 6}
Il y a $5+7+8+10=30$ élèves. La médiane est donc la moyenne des 15\ieme{} et 16\ieme{} valeurs de la série ordonnée.

Les notes 7 occupent les rangs 1 à 5, les notes 10 les rangs 6 à 12, les notes 12 les rangs 13 à 20. Les 15\ieme{} et 16\ieme{} valeurs valent donc toutes les deux 12. La médiane est donc $12$.

La bonne réponse est \textbf{a.}

\subsubsection*{Question 7}
Sur la courbe, la fonction est positive ou nulle lorsque la courbe est située sur ou au-dessus de l'axe des abscisses.

On lit que la courbe coupe l'axe des abscisses en $x=-3$ et en $x=-1$. Elle est au-dessus de l'axe sur $[-4\,;\,-3]$ puis sur $[-1\,;\,3]$. L'ensemble des solutions est donc : $[-4\,;\,-3]\cup[-1\,;\,3]$.

La bonne réponse est \textbf{c.}

\subsubsection*{Question 8}
La proportion d'élèves pratiquant un sport est : $\dfrac{150}{500}=0{,}3=30\,\%$.

La bonne réponse est \textbf{c.}




\begin{center}
\end{center}
